\subsubsection{Tauri აპლიკაცია (desktop კლიენტი)}

Tauri აპლიკაციის Rust ბექენდი ის ადგილია, სადაც ყველა ქვესისტემა
ერთმანეთს ხვდება (იხ. დიაგრამა~\ref{diag:layers}): დოკუმენტის აქტორი,
WebSocket კავშირი, სესიის მდგომარეობის მანქანა, პროცესების მართვა, ნაგვის
შემგროვებელი, ფაილების შენახვა, უფლებების სისტემა და მეტრიკები.

\paragraph{დოკუმენტის აქტორი.}
ყველაზე მნიშვნელოვანი დიზაინის გადაწყვეტილება დოკუმენტთან წვდომის
ორგანიზებაა. გულუბრყვილო მიდგომა იქნებოდა \texttt{Mutex<Document>} -
საზიარო ობიექტი საკეტით, რომელსაც ყველა ნაკადი ეპოტინება. ამის ნაცვლად
\emph{აქტორის} ნიმუში გამოვიყენეთ: დოკუმენტს ფლობს ერთადერთი tokio
ამოცანა (\texttt{DocActor}), რომელიც უსასრულო ციკლში კითხულობს
შეტყობინებების არხს (\texttt{mpsc}). ნებისმიერი მოქმედება, ლოკალური
ჩასმა, დაშორებული ოპერაცია, სნეფშოტის აღება, ფაილში ჩასაწერი ტექსტის
წაკითხვა, არხში იგზავნება \texttt{DocOp} შეტყობინებად, პასუხი კი
ერთჯერადი (oneshot) არხით ბრუნდება. ამ მიდგომას რამდენიმე უპირატესობა
აქვს: საკეტები და deadlock-ის რისკი ქრება, ოპერაციები ბუნებრივად
სერიალიზდება (რაც CRDT-ისთვის ისედაც აუცილებელია), ხოლო ახალი ოპერაციის
დამატება მხოლოდ ახალი \texttt{DocOp} ვარიანტის და მისი დამმუშავებლის
დაწერას ნიშნავს. Rust-ის ownership სისტემა აქ ჩვენს სასარგებლოდ მუშაობს -
დოკუმენტზე მიმართვა აქტორის გარეთ უბრალოდ ვერ „გაიპარება“.

\paragraph{IPC, ბექენდიდან ფრონტენდამდე.}
როდესაც აქტორი დაშორებულ ცვლილებას აინტეგრირებს, ფრონტენდს Tauri ივენთს
უგზავნის: \texttt{crdt://remote-change} (პოზიცია და შიგთავსი),
\texttt{crdt://snapshot-applied} (დოკუმენტი მთლიანად შეიცვალა, მაგ.
სესიაში შესვლისას) ან \texttt{crdt://document-reset}. საპირისპირო
მიმართულებით ფრონტენდი \texttt{\#[tauri::command]} ბრძანებებს იძახებს:
\texttt{insert}, \texttt{delete}, \texttt{replace}, ფაილის ოპერაციები,
სესიის მართვა და ა.შ. საინტერესო ნიუანსი: webview-სა და ბექენდს შორის
შეტყობინებები ასინქრონულია, ამიტომ შესაძლებელია, რომ მომხმარებელმა
პოზიციაზე დააჭიროს მაშინ, როცა ბექენდში უკვე ჩამოვიდა (მაგრამ რედაქტორში
ჯერ არ ასახულა) სხვისი ცვლილება. ამის გამო ყოველი ლოკალური ბრძანება თან
ატარებს \texttt{base\_seq}-ს, ბოლო დაშორებული ცვლილების ნომერს, რომელიც
რედაქტორს ჰქონდა ნანახი ბრძანების გაგზავნისას. აქტორი ინახავს ბოლო
დაშორებული ოპერაციების მცირე ჟურნალს და შემოსული პოზიციის მათ მიმართ
ტრანსფორმაციას აკეთებს, პატარა, OT-ის სულისკვეთების ხიდი webview-სა და
ბექენდის რეპლიკას შორის, რომელმაც სინქრონიზაციის რთულად დასაჭერი ბაგები
აღმოფხვრა.

\paragraph{პროცესების მართვა.}
სესიის დაწყებისას \texttt{process\_coordinator} უშვებს გეითვეის sidecar-ს,
სტანდარტული გამოსავლიდან კითხულობს \texttt{\{"port": N\}} შეტყობინებას,
შემდეგ (საუკეთესო ძალისხმევით) უშვებს \texttt{cloudflared}-ს და ლოგებიდან
იჭერს გაცემულ საჯარო ბმულს. გეითვეისთან ავტორიზაციისთვის აპლიკაცია ყოველ
სესიაზე შემთხვევით bearer ტოკენს აგენერირებს და პროცესს გარემოს ცვლადით
გადასცემს, ამ ტოკენით ხდება ოთახის შექმნა და სესიის დასრულება (თავად
WebSocket შეერთება ტოკენს არ ითხოვს, რადგან სტუმრებს მისი ცოდნა არ
სჭირდებათ და არც უნდა სჭირდებოდეთ). ცალკე ყურადღება დაეთმო პროცესების
სიცოცხლის ციკლს: ფანჯრის დახურვისას აპლიკაცია სინქრონულად შლის ოთახს და
კლავს ორივე sidecar-ს, რომ მომხმარებელს „ობოლი“ სერვერი არ დარჩეს
ფონურად გაშვებული.

\paragraph{ასინქრონულობა და WebSocket მართვა.}
მთელი ბექენდი tokio-ს ასინქრონულ runtime-ზე დგას. WebSocket კავშირს სამი
დამოუკიდებელი tokio ამოცანა ემსახურება: \texttt{write\_loop} (გამავალი
ფრეიმების რიგი), \texttt{receive\_loop} (შემომავალი ბაიტების მიღება და
ტიპის ამოცნობა) და \texttt{process\_loop} (მიღებული ოპერაციების
გადაცემა დოკუმენტის აქტორისთვის და სხვა ქვესისტემებისთვის). ამოცანების ეს
გამიჯვნა ორ რამეს გვაძლევს: ნელი ქსელი ვერასდროს ბლოკავს დოკუმენტის
დამუშავებას, ხოლო კავშირის გაწყვეტისას საკმარისია სამივე ამოცანის
გაუქმება და მდგომარეობის გაწმენდა ერთ ადგილას (disconnect handler).

\paragraph{სესიის მდგომარეობის მანქანა.}
სესიის მდგომარეობა მკაცრი სასრული ავტომატია (FSM):
\texttt{Undecided} $\rightarrow$ \texttt{Starting} $\rightarrow$
\texttt{Host}/\texttt{Guest} $\rightarrow$ \texttt{Undecided}. ყველა
გადასვლა ვალიდირდება და უკანონო გადასვლა შეცდომას აბრუნებს. სესიის
წამოწყება მრავალსაფეხურიანი ასინქრონული პროცესია (პროცესების გაშვება,
ოთახის შექმნა, WebSocket კავშირი) და ნებისმიერ საფეხურზე შეიძლება
ჩავარდეს, ამიტომ \texttt{Starting} მდგომარეობას RAII „მცველი“ (guard)
იცავს: თუ ფუნქცია შეცდომით დასრულდა და მცველი ისე განადგურდა, რომ სესია
არ დასრულებულა, მდგომარეობა ავტომატურად ბრუნდება \texttt{Undecided}-ზე.
ასე გამოვრიცხეთ „გაჭედილი“ შუალედური მდგომარეობები ხელით rollback-ების
წერის გარეშე.

\paragraph{უფლებების მართვა.}
ჩაწერის უფლებების ავტორიტეტი გეითვეია: ის პირველ შემოერთებულს ჰოსტად
თვლის, ყოველ კლიენტს ჩაწერის დროშას უნახავს და წაკითხვის რეჟიმში მყოფი
კლიენტის ოპერაციებს უბრალოდ აგდებს (კლიენტის კოდის მოდიფიცირებული ვერსიაც
რომ შეეცადოს წესის გვერდის ავლას, გეითვეი მის ოპერაციებს არ გაავრცელებს).
ჰოსტი უფლების შეცვლას სპეციალური ფრეიმით (\texttt{0x07}) ითხოვს, გეითვეი
ამოწმებს, რომ მოთხოვნა ნამდვილად ჰოსტისგან მოდის, ცვლის დროშას და იმავე
ფრეიმს ყველას ავტორიტეტულ ექოდ უგზავნის. კლიენტის მხარეს
\texttt{roster} მოდული აწარმოებს მონაწილეთა სიას (სახელები, ჰოსტობა,
უფლებები), ხოლო სტუმრის საკუთარი უფლება \texttt{AppRole::Guest}-შივე
ინახება და \texttt{insert}/\texttt{delete}/\texttt{replace} ბრძანებებს
კარიბჭესავით კეტავს. ინტერფეისში ეს Monaco-ს \texttt{readOnly} რეჟიმად
აისახება. აქ ერთი ხრიკი დაგვჭირდა: Monaco წაკითხვის რეჟიმში პროგრამულ
ცვლილებებსაც ჩუმად აგდებს, ამიტომ დაშორებული ცვლილების ჩასმისას დროშას
წამიერად ვხსნით და ისევ ვაბრუნებთ.

\paragraph{ფაილების შენახვა და გახსნა.}
დოკუმენტი ორ ფორმატში ინახება, გაფართოების მიხედვით: ჩვეულებრივი ტექსტი
(ნებისმიერი გაფართოება) ან საკუთარი \texttt{.pcdoc} ფორმატი, \texttt{PCDC}
მაგიური ბაიტები, ვერსია და ბინარულად სერიალიზებული სრული CRDT სნეფშოტი,
ისტორიითა და იდენტიფიკატორებით. გახსნისას ფორმატს ვცნობთ არა
გაფართოებით, არამედ მაგიური ბაიტებით. ჩაწერა ყოველთვის ატომურია, ჯერ
დროებით ფაილში, შემდეგ გადარქმევით, რომ შენახვისას შეწყვეტამ (გათიშულმა
დენმა, ჩავარდნილმა პროცესმა) ნახევრად ჩაწერილი ფაილი არ დაგვიტოვოს.
ცალკე პრობლემა აღმოჩნდა ხაზის დაბოლოებები: Windows-ის ფაილები CRLF-ს
იყენებს, Monaco და CRDT კი LF-ს ელოდება. გახსნისას ტექსტს LF-ზე
ვანორმალიზებთ, ვიმახსოვრებთ, იყო თუ არა ფაილი CRLF-იანი, და შენახვისას
ისევ ვაბრუნებთ, ასე Windows-ის ფაილი ბაიტ-ბაიტ იდენტური ბრუნდება და
კროსპლატფორმულ სესიაშიც პოზიციები აღარ „ცურავს“ (ეს ბაგი რეალურად
გვქონდა და ქვემოთ, პრობლემების ნაწილში, ცალკე ვახსენებთ). ბოლოს,
აპლიკაცია ბოლო გახსნილი ფაილების სიას (MRU) აწარმოებს, სრულფასოვანი
ფაილ-მენეჯერის გარეშე, უბრალოდ სწრაფი წვდომისთვის.

\paragraph{მეტრიკების სისტემა.}
დიაგნოსტიკისთვის აპლიკაციას ჩაშენებული მეტრიკების პანელი აქვს. გეითვეი
ატომურ მთვლელებს აწარმოებს, აქტიური ოთახები და კლიენტები, გადაგზავნილი
ფრეიმების რაოდენობა და მოცულობა, სნეფშოტ-სინქრონიზაციების
წარმატება/ჩავარდნა, ნელი კლიენტების გათიშვები, და მათ HTTP ენდპოინტზე
აქვეყნებს. ჰოსტის აპლიკაციაში ცალკე tokio ამოცანა
(\texttt{process\_metrics\_aggregator}) პერიოდულად კითხულობს გეითვეის ამ
ენდპოინტს და \texttt{cloudflared}-ის საკუთარ მეტრიკებს, აერთიანებს მათ და
ივენთით ფრონტენდს აწვდის, სადაც მომხმარებელი მათ პანელში ხედავს. ეს
სესიის გამართვისას ძალიან გამოგვადგა, მაშინვე ჩანს, გადის თუ არა
ტრაფიკი და ვინ არის შეერთებული.
